\section{Regularizing the actual tower and excluding return defects}
\label{sec:actual-return-defects}
A function of an actual return state lifts to the collision tower with
an exactly localized defect. That lift need not have bounded variation.
We construct a regular approximation with a controlled cost before
using Section~\ref{sec:near-origin-defects}. The construction concerns
functions and first distributional derivatives, not anisotropic
vectors or inverse-coarea densities.

\subsection{Finite collision records in the initial coordinates}
\begin{lemma}[A quantitative finite-record variation bound]
\label{lem:finite-record-variation}
There is $C\ge1$ such that, for every $L\ge1$ and $R\in I$, the
following bounded functions have variation norm at most
\begin{equation}\label{eq:finite-record-variation-budget}
 B_L=\exp\{C L\log(2+L)\}:
\end{equation}
all components of
$\chi\circ T_R^j$, $h_R\circ T_R^j$, and $\eta_R\circ T_R^j$,
$|j|\le L$, where $\chi(\alpha,p)=(\cos\alpha,\sin\alpha,p)$.
Their supremum norms are bounded independently of $L$. Null singular
sets are assigned any fixed bounded semialgebraic convention.
Moreover, if $g$ is smooth on the reflected collision cylinder, then
\begin{equation}\label{eq:finite-record-smooth-composition}
 \|g\circ T_R^j\|_{\mathrm{BV}}
              \le C\|g\|_{C^1(\alpha,p)}B_L\qquad(|j|\le L).
\end{equation}
The constant in $B_L$ is independent of the coefficients specifying
the moving section rectangles.
\end{lemma}
\begin{proof}
We supply the finite-complexity argument, since a one-collision BV
bound does not alone imply this lemma. A fixed finite set of disk
labels contains every next collision, by the uniform horizon. Use
normal vectors $n_j\in\mathbb S^1$, outgoing velocities
$v_j\in\mathbb S^1$, and flight lengths $\tau_j$ as auxiliary
variables. For a fixed label word $d_j$, each step is given by
\begin{equation}\label{eq:polynomial-flight-graph}
 Rn_j+\tau_jv_j-d_j=Rn_{j+1},\qquad
 v_{j+1}=v_j-2(v_j\cdot n_{j+1})n_{j+1},\qquad
 p_j=v_j\cdot Jn_j,
\end{equation}
where $J$ rotates vectors by $\pi/2$. Require $\tau_j>0$,
$v_j\cdot n_j>0$, and $v_j\cdot n_{j+1}<0$.
Require also that the open flight not enter any other candidate disk.
These last conditions are finite Boolean combinations of polynomial
sign conditions of bounded degree: for each center $d$, minimize
$|Rn_j+t v_j-d|^2$ over $0\le t\le\tau_j$.
The minimizing value occurs at an endpoint or at
$t=(d-Rn_j)\cdot v_j$; splitting into these three cases gives the
required inequalities without a quantified time variable. Endpoints
and tangencies are included in the declared singular convention.
Thus the graph records the actual first collision, not merely some
root of the flight equation.

On each bounded rational angular chart, the initial normal and outgoing
velocity also have such a graph, using a positive square-root variable
for $\sqrt{1-p^2}$. A word of length at most $2L+2$ therefore uses
at most $C_2L$ variables and $C_2L$ polynomials, of degree bounded by
a fixed $d$. All these variables can be confined to a bounded ball.
The same representation gives negative iterates by fixing the terminal
rather than the initial state; the regular billiard map is invertible.
Every lifted fiber over a regular initial state is unique.

Section membership is equally finite. Each angular interval has length
less than $\pi$ and is described by two oriented half-plane tests on
$n_j$, together with the two inequalities for $p_j$. The moving angular
centers enter only as real coefficients. Taking the finite union of
rectangles gives $\eta_R$; adjoining its two possible values also gives
the graph of $h_R=f_R-\bar G_R\eta_R$. All degree and variable counts
remain as above. There are at most $C_3^L$ choices of label word and
of these finite membership alternatives.

Here is the topological bound we use. The finite sign-condition bound
recorded in \cite[Section 3.2, equation (3.3)]{BPR}, with the zeroth
Betti number and ambient dimension $k$, bounds the total number of
connected components by
\[
 d(2d-1)^{k-1}\sum_{i=0}^{k}\binom{s}{i}4^i
                       \le (C(1+s)d)^{Ck}.
\]
The bound applies uniformly in the polynomial coefficients and to any
Boolean union of the sign conditions. We use this finite inequality,
not an asymptotic formula whose dimension is held fixed.
Fix the second collision coordinate and a real level $a$. Add the
condition that a chosen scalar output $u$ exceed $a$. The lifted
semialgebraic set has the preceding complexity. Its projection onto
the remaining coordinate has at most as many connected components:
every connected component projects to an interval or a point. Summing
over the at most $C_3^L$ alternatives bounds the number of components
of the slice superlevel set by $\exp\{C L\log(2+L)\}$.
The same reasoning applies with the two input coordinates interchanged.

If $|u|\le M_0$, the perimeter of a one-dimensional superlevel set
with at most $b$ components is at most $2b$ in the interior. The
one-dimensional coarea formula therefore bounds the variation of
that slice by $4M_0b$. Integrating both slice bounds gives the two
first distributional derivative bounds, including jumps. The fixed
angular chart cover has bounded derivatives and inverse derivatives;
its seams add only uniformly bounded boundary terms. This proves
\eqref{eq:finite-record-variation-budget}, after increasing $C$.
No inverse derivative near grazing has been bounded pointwise.

Finally extend $g$ from the embedded cylinder
$\chi(M)\subset\R^3$ to a neighborhood, with Lipschitz norm at most
$C\|g\|_{C^1(\alpha,p)}$, by a fixed tubular chart and cutoff. The
BV chain inequality applied to the three components of
$\chi\circ T_R^j$ proves \eqref{eq:finite-record-smooth-composition}.
This does not assert a BV composition bound for an arbitrary nonsmooth
$g$; only the smoothed functions below require this estimate.
\end{proof}

\subsection{Exact lifting and its regular approximation}
Write $Y=Y_R^*$, $F=F_R^*$, $r_R^*$ for its actual return time and
$g_R=G_R-\bar G_R$. For almost every collision state let its age since
the most recent section visit be $j$. The corresponding disjoint levels
are
\begin{equation}\label{eq:actual-age-levels}
 D_{j,R}=T_R^j\{y\in Y:r_R^*(y)>j\},\qquad j\ge0.
\end{equation}
They partition a full-measure set, by recurrence and invertibility.
For a measurable section function $f$, define its exact phase lift by
\begin{equation}\label{eq:exact-phase-tower-lift}
 (\mathsf L_{R,z}f)(T_R^jy)
       =e^{iz\cdot S_{j,R}h_R(y)}f(y),\qquad 0\le j<r_R^*(y).
\end{equation}
This is a function lift, not an induced transfer operator.

\begin{lemma}[The defect is supported at the true tower tops]
\label{lem:exact-tower-defect}
For $1\le p<\infty$, whenever the terms are finite,
\begin{align}
 \|\mathsf L_{R,z}f\|_{L^p(\nu)}^p
       &=c_*\int_Y r_R^*|f|^p\dd\nu_R^*,
                      \label{eq:tower-lift-norm}\\
 \|\mathcal D_{R,z,0}(\mathsf L_{R,z}f)\|_{L^p(\nu)}^p
       &=c_*\int_Y|f\circ F-e^{iz\cdot g_R}f|^p\dd\nu_R^*.
                      \label{eq:tower-exact-defect}
\end{align}
The analogous supremum bound is
$\|\mathsf L_{R,z}f\|_\infty\le\|f\|_\infty$.
\end{lemma}
\begin{proof}
Invariance and the disjoint levels give, for every nonnegative $b$,
\[
 \int_M b\dd\nu
   =c_*\int_Y\sum_{j=0}^{r_R^*-1}b(T_R^jy)\dd\nu_R^*(y).
\]
Apply this first to the modulus of the lift. At every nontop level
its next value is exactly $e^{iz\cdot h_R}$ times its current value.
At the last level the defect is
\[
 f(Fy)-e^{iz\cdot S_{r_R^*,R}h_R(y)}f(y)
               =f(Fy)-e^{iz\cdot g_R(y)}f(y),
\]
using the one-return compensation identity. This proves
\eqref{eq:tower-exact-defect}. Only one top occurs per return; its
measure is integrated against $c_*\nu_R^*$, not against a
return-time-biased probability.
\end{proof}

For a section function $f$ let $f^0$ be its zero extension to $M$,
$M_f=\|f\|_\infty$ and $V_f=\|f^0\|_{\mathrm{BV}}$.
Set $f_\epsilon=\mathsf S_\epsilon f^0$ and define
\begin{equation}\label{eq:finite-age-approximation}
 Q_{L,\epsilon}(x)=\sum_{j=0}^{L-1}\one_{D_{j,R}}(x)
 e^{iz\cdot\sum_{i=1}^j h_R(T_R^{-i}x)}
                              f_\epsilon(T_R^{-j}x).
\end{equation}
The sum for $j=0$ is empty. The approximant is zero on ages at least
$L$; the original lift is not truncated in any theorem statement.

\begin{proposition}[A uniform cost for approximating the lift]
\label{prop:tower-bv-approximation}
For $L\ge2$ and $0<\epsilon<1/2$, uniformly in $R,z,f$,
\begin{align}
 \|Q_{L,\epsilon}\|_\infty&\le M_f,
                                      \label{eq:tower-sup-bound}\\
 \|\mathsf L_{R,z}f-Q_{L,\epsilon}\|_1
        &\le C(L\epsilon V_f+M_fe^{-cL}),
                                      \label{eq:tower-L1-approximation}\\
 \|\mathsf L_{R,z}f-Q_{L,\epsilon}\|_2^2
        &\le C(M_fL\epsilon V_f+M_f^2e^{-cL}),
                                      \label{eq:tower-L2-approximation}\\
 \|Q_{L,\epsilon}\|_{\mathrm{BV}}
        &\le M_f\exp\{C L\log(2+L)\}(1+\epsilon^{-1}+|z|).
                                      \label{eq:tower-BV-budget}
\end{align}
\end{proposition}
\begin{proof}
The level indicators are disjoint, proving the supremum bound.
On the first $L$ levels, invariance bounds the total $L^1$ change
of the initial function by
$L\|f^0-f_\epsilon\|_1\le CL\epsilon V_f$.
The remaining age probability is exactly
\[
 \nu\{\text{age}\ge L\}
       =c_*\int_Y(r_R^*-L)_+\dd\nu_R^*\le Ce^{-cL},
\]
by Theorem~\ref{thm:exponential-return}. This proves the $L^1$
bound. Both functions have supremum at most $M_f$; multiplying the
$L^1$ bound by $2M_f$ gives the $L^2$ bound.

For variation, the exact level indicator is
\[
 \one_{D_{j,R}}=(\eta_R\circ T_R^{-j})
                    \prod_{i=0}^{j-1}(1-\eta_R\circ T_R^{-i}).
\]
The finite-record lemma and the BV product inequality bound its
variation by $C(j+1)B_L$. The same lemma bounds the phase sum
variation by $CjB_L$. The exponential is a Lipschitz function of this
real vector sum, with Lipschitz constant at most $|z|$.
Moreover $\|f_\epsilon\|_{C^1}\le C M_f\epsilon^{-1}$, so
\eqref{eq:finite-record-smooth-composition} applies to its pullback.
Apply the BV product inequality to each summand and then sum over
$j<L$. Absorb the resulting fixed powers of $L$ into $B_L$ by
increasing its constant. This proves \eqref{eq:tower-BV-budget}.
Neither the levels nor their images have been declared smooth; their
indicator jumps are included in the BV estimates.
\end{proof}

\subsection{Defects for the actual unbounded return record}
For $H\ge2$, $0<r<1/2$ define
\begin{equation}\label{eq:return-log-budget}
 \Lambda=1+\log(H/r),\qquad \mathcal K(H,r)=\Lambda\log(2+\Lambda).
\end{equation}
The letter $\mathcal K$ here denotes a scalar budget, not a compact
frequency set as in Section~\ref{sec:quantitative-phase-defects}.

\begin{theorem}[Uniform return-phase coercivity near the origin]
\label{thm:actual-return-circle-defect}
There exist $a_2,c_2,r_2>0$ such that, for $H\ge2$,
$0<r=|z|\le r_2$, and $r\mathcal K(H,r)\le a_2$,
\begin{equation}\label{eq:actual-return-circle-defect}
 |q|=1\text{ on }Y_R^*,\quad\|q^0\|_{\mathrm{BV}}\le H
 \quad\Longrightarrow\quad
 \|q\circ F_R^*-e^{iz\cdot(G_R-\bar G_R)}q\|_{L^1(\nu_R^*)}
                         \ge c_2r^2.
\end{equation}
All constants are independent of $R$, and no induced mixing hypothesis
is used.
\end{theorem}
\begin{proof}
The exact lift $Q=\mathsf L_{R,z}q$ has modulus one. Given
$\eta=b r^2$ with $b>0$ fixed and small, choose
$L\le C(1+|\log\eta|)$ and
$\epsilon=\eta/(C L H)$ so that
\eqref{eq:tower-L1-approximation} gives
$\|Q-Q_{L,\epsilon}\|_1\le\eta$.
The circle truncation in Lemma~\ref{lem:bv-circle-truncation}, applied
to $Q_{L,\epsilon}$, gives a phase $\widetilde q$ satisfying
\[
 \|\widetilde q-Q\|_1\le4\eta,\qquad
 1+\log(2+\|\widetilde q\|_{\mathrm{BV}})
                                     \le C\mathcal K(H,r).
\]
For the first bound, use the pointwise inequality from that lemma,
$|\widetilde q-Q_{L,\epsilon}|\le3\big||Q_{L,\epsilon}|-1\big|
\le3|Q_{L,\epsilon}-Q|$. The second follows from
\eqref{eq:tower-BV-budget}; the constants include the chosen $b$.

After decreasing $a_2,r_2$, Theorem~\ref{thm:near-origin-circle-defect}
applies to $\widetilde q$. Changing a function by $4\eta$ in $L^1$
changes its collision defect by at most $8\eta$.
Equation~\eqref{eq:tower-exact-defect} therefore gives
\[
 c_0r^2\le c_*\|q\circ F_R^*-e^{iz\cdot g_R}q\|_1+8br^2.
\]
Choose $b<c_0/16$. This proves the assertion with a uniform $c_2>0$.
In particular the regularity of the exact lift was not assumed.
\end{proof}

\begin{theorem}[Uniform coercivity for actual return functions]
\label{thm:actual-return-vector-defect}
There exist $a_3,c_3,r_3>0$ such that
$0<r=|z|\le r_3$, $H\ge2$, and $r\mathcal K(H,r)\le a_3$ imply
\begin{align}
 &\|f\|_{L^2(\nu_R^*)}=1,\quad M_f+V_f\le H
 \quad\Longrightarrow\notag\\
 &\qquad
 \|f\circ F_R^*-e^{iz\cdot(G_R-\bar G_R)}f\|_{L^2(\nu_R^*)}
                     \ge\frac{c_3r^2}{\mathcal K(H,r)}.
                         \label{eq:actual-return-vector-defect}
\end{align}
The section function may be complex and may vanish. The conclusion
is uniform in the radius and in the direction of $z$.
\end{theorem}
\begin{proof}
Let $Q=\mathsf L_{R,z}f$ and $S=\|Q\|_2$. By
\eqref{eq:tower-lift-norm}, $\sqrt{c_*}\le S\le H$; the upper bound
also follows from $\|Q\|_\infty\le H$ and $\nu(M)=1$.
If $e$ denotes the defect on the left of
\eqref{eq:actual-return-vector-defect}, then
\begin{equation}\label{eq:normalized-tower-defect}
 \|\mathcal D_{R,z,0}(Q/S)\|_2=\sqrt{c_*}\,e/S\le e.
\end{equation}

Choose an accuracy
$\eta=b r^2/\mathcal K(H,r)$, where $0<b<1$ will be fixed.
The bounds in Proposition~\ref{prop:tower-bv-approximation} allow
$L\le C[\Lambda+|\log b|]$ and
$\epsilon\ge c\eta^2/(H^2L)$ with
$\|Q-Q_{L,\epsilon}\|_2\le\sqrt{c_*}\eta$.
For example choose $L$ to make the exponential term at most
$c_*\eta^2/2$, then choose $\epsilon$ equal to a sufficiently small
fixed multiple of $\eta^2/(H^2L)$. Both choices are uniform in $R$.
Normalize $Q_{L,\epsilon}/S$ to an $L^2$ unit function $v$.
For $\eta<1/4$ this is well defined and
\[
 \|v-Q/S\|_2\le2\eta,\qquad
 1+\log(2+\|v\|_\infty+\|v\|_{\mathrm{BV}})+|\log r|
       \le C(1+|\log b|)^2\mathcal K(H,r).
\]
The last bound is \eqref{eq:tower-BV-budget}, using the lower bound
on $S$, and $\log\mathcal K(H,r)\le C\Lambda$.
The displayed dependence on $b$ prevents a circular choice of accuracy.

Choose $b$ so small that
$4b<c_1/[2C(1+|\log b|)^2]$; such a choice exists because
$b(1+|\log b|)^2\to0$. Fix this $b$ once and for all. Decrease
$a_3,r_3$ so that Theorem~\ref{thm:near-origin-vector-defect} applies
to $v$ under the preceding budget. That theorem and
\eqref{eq:normalized-tower-defect} give
\[
 e+4\eta\ge\|\mathcal D_{R,z,0}v\|_2
   \ge\frac{c_1r^2}{C(1+|\log b|)^2\mathcal K(H,r)}.
\]
The selected accuracy absorbs $4\eta$ and proves
\eqref{eq:actual-return-vector-defect}. No lower bound for $|f|$
and no mixing rate for $F_R^*$ enter the proof.
\end{proof}

\begin{corollary}[The intervening annulus with polynomial budgets]
\label{cor:actual-annulus-defect}
Fix $H_0\ge2$ and $\zeta\ge0$. For all sufficiently large $n$,
uniformly in $R$ and in physical frequencies
\begin{equation}\label{eq:actual-annulus-frequencies}
 2n^{-99/200}\le|z|\le n^{-2/5},\qquad
        2n^{1/200}\le|\sqrt n\,z|\le n^{1/10},
\end{equation}
unit section phases with zero-extension variation at most $H_0n^\zeta$
have defect at least $c|z|^2$. Normalized complex section functions
with $M_f+V_f\le H_0n^\zeta$ have defect at least
\begin{equation}\label{eq:actual-annulus-vector-rate}
 \frac{c|z|^2}{\log(2+n)\log\log(e^e+n)}.
\end{equation}
In particular this last lower bound is at least
$c n^{-99/100}/[\log(2+n)\log\log(e^e+n)]$ on the whole annulus.
\end{corollary}
\begin{proof}
On this annulus $\Lambda(H_0n^\zeta,|z|)\le C\log(2+n)$, so
$\mathcal K\le C\log(2+n)\log\log(e^e+n)$.
The largest possible value of $|z|\mathcal K$ is bounded by
$C n^{-2/5}\log(2+n)\log\log(e^e+n)$, which tends to zero.
Thus both preceding theorems apply. Squaring the lower frequency
endpoint gives the final power. The interval is nonempty for all
sufficiently large $n$, since its endpoint ratio is
$n^{19/200}/2\to\infty$.
\end{proof}

\paragraph{What the annulus estimate supplies.}
The norm in \eqref{eq:actual-annulus-vector-rate} is a one-step
cohomological defect for the \emph{actual unbounded return record},
with the normalization and first-variation budget stated above.
The lifted defect and the regularization error have both been proved,
rather than assumed to come from an induced spectral construction.
No conclusion is made here for arbitrary eigenvalues of the induced
operator: the multiplier in the theorem is precisely
$e^{iz\cdot(G_R-\bar G_R)}$.

This estimate is a quantitative input on the intervening annulus, not
the complementary Fourier integral itself. To deduce decay of powers
from it on a distribution space still requires a construction of that
space, a function reconstruction with the displayed budgets, and a
bound relating the operator defect to the present physical $L^2$
defect. Nor does the finite-record BV lemma estimate any second
derivative of a raw pushforward density. The growing roof region,
the complete critical/singular extraction, and the weighted exact-event
comparison remain distinct requirements of Theorem~\ref{thm:LLT}.
All conditioning indicators in the inherited central and moment
theorems are unchanged.
